\scanpage{135}
\unsure{1|I} \uline{Lemme} Soit $S$ un préschéma loc noeth, et
pour \deleted{tout} $s\in S$, soit $\deleted{\unsure{X_s}}$ $X(s)$ un préschéma
quasi fini \uline{étale} \uline{séparé} sur $\mathrm{Spec}(\mathcal O_s)$, et pour
$t\in S$, $s\in\overline t$, soit $\varphi_{st}:X(t)\unsure{-|\to}X(s)$ rendant
\diagram{\begin{tikzcd}[ampersand replacement=\&,row sep=small]
X(s)\arrow[d] \& X(t)\arrow[l]\arrow[d]\\
\mathrm{Spec}(\mathcal O_s) \& \mathrm{Spec}(\mathcal O_t)\arrow[l]
\end{tikzcd}}
cartésien le diagramme ci-contre. On suppose qu'on a transitivité
$\varphi_{st}\varphi_{tu}=\varphi_{su}$.
% margin note written sideways, linked to the end of the transitivity condition
\inserted{[On suppose que pour tt $s\in S$, il existe une \unsure{réalisation} étale \deleted{\illegible{1}}
$\mathfrak X$ de $X(s)$ sur un voisinage ouvert $U$ de $s$, \deleted{\unsure{tel que pour tt}} et pour tt $t\in U$,
un \unsure{hom} $\deleted{\unsure{\mathfrak X_s}}\mathfrak X\deleted{\unsure{s}}\times_U\mathrm{Spec}(\mathcal O_t)\to X(t)$, ces hom étant compatibles
avec les \unsure{\ldots}]}
% margin note written sideways, diagram alongside the inserted condition
\diagram{\begin{tikzcd}[ampersand replacement=\&,column sep=small,row sep=small]
X(s)\arrow[r]\arrow[d] \& \mathfrak X\arrow[d] \& \mathfrak X\times_U\mathrm{Spec}(\mathcal O_t)\arrow[l]\arrow[r]\arrow[d] \& X(t)\\
\mathrm{Spec}(\mathcal O_s)\arrow[r] \& U \& \mathrm{Spec}(\mathcal O_t)\arrow[l] \&
\end{tikzcd}}

\deleted{\unsure{des}} Alors il existe un préschéma
étale \inserted{quasi fini séparé} sur $S$, soit $X$, et des \inserted{$\mathcal O_s$-}isom
\[
\varphi_s:X\times_S\mathrm{Spec}(\mathcal O_s)\unsure{-|\to|\overset{\sim}{\longrightarrow}}X(s)
\]
rendant commutatifs les diagrammes qu'on pense.
$X$ est unique à isom unique près.

Ce dernier \deleted{fait} résulte d'un
énoncé plus général sur la catégorie des
préschémas de type fini sur $S$, qui donne
un foncteur \unsure{pleinement} fidèle. Il nous permet
de nous ramener au local et supposer $S$
affine, donc noeth. On procède par
récurrence noethérienne, en supposant l'énoncé
prouvé pour les sous-préschémas \unsure{fermés} de
$S\ne S$.
% archive sheet number: 134
